% CAHIER DES CHARGES
\chapter*{CAHIER DES CHARGES}
\addcontentsline{toc}{chapter}{Cahier des charges}
\thispagestyle{plain}

\textbf{SOUS-PROJET 2 : Planification et optimisation des trajectoires}\\
\textbf{(Robotique -- STIC)}

\vspace{0.5cm}
\textbf{Objectif et problématique :}

Concevoir et optimiser les trajectoires d'un robot à câbles afin de couvrir efficacement une parcelle agricole tout en minimisant le temps de parcours et la consommation énergétique.

Comment définir des trajectoires optimales pour un robot à câbles en tenant compte des contraintes mécaniques, de la visibilité des plantes et du compromis entre couverture, précision et efficacité énergétique ?

\vspace{0.35cm}
\textbf{Méthodologie}

\begin{enumerate}[leftmargin=*,itemsep=2pt]
  \item Modélisation de l'environnement (parcelle, points d'intérêt)
  \item Génération de trajectoires (balayage, points d'intérêt)
  \item Prise en compte des contraintes (limites mécaniques, visibilité)
  \item Optimisation (temps, énergie, couverture)
  \item Simulation (implémentation Python/MATLAB, validation)
\end{enumerate}

\vspace{0.25cm}
\textbf{Résultats attendus :}
\begin{itemize}[leftmargin=*,itemsep=2pt]
  \item Trajectoires optimisées
  \item Simulation fonctionnelle
  \item Comparaison des méthodes
\end{itemize}
